% ========================================================
%  《习近平新时代中国特色社会主义思想概论》全真模拟试卷（含参考答案）- 第二套
%  覆盖范围：导论、第一章至第十二章核心考点
% ========================================================
\documentclass[12pt, a4paper]{article}

% ---- 中文支持 ----
\usepackage[UTF8]{ctex}

% ---- 数学 ----
\usepackage{amsmath, amssymb, amsthm}
\usepackage{mathtools}

% ---- 页面布局 ----
\usepackage[top=2.5cm, bottom=2.5cm, left=2.8cm, right=2.8cm]{geometry}

% ---- 表格 ----
\usepackage{booktabs, array, multirow, makecell}
\usepackage{tabularx}

% ---- 页眉页脚 ----
\usepackage{fancyhdr}
\usepackage{lastpage}
\pagestyle{fancy}
\fancyhf{}
\renewcommand{\headrulewidth}{0.6pt}
\lhead{\small 习近平新时代中国特色社会主义思想概论·GLM·B卷}
\rhead{\small 共 \pageref{LastPage} 页}
\cfoot{\small 第 \thepage\ 页}

% ---- 计数器与样式 ----
\usepackage{enumitem}
\usepackage{xcolor}
\usepackage{tcolorbox}
\tcbuselibrary{skins, breakable}

% ---- 填空线 ----
\newcommand{\blank}[1]{\underline{\hspace{#1}}}

% ---- 答案盒子 ----
\newtcolorbox{answerbox}[1][]{
  breakable,
  colback=gray!6,
  colframe=gray!50,
  fonttitle=\bfseries\small,
  title=参考答案,
  #1
}

% ---- 题号格式 ----
\newcounter{bigq}
\newcommand{\BigQ}[1]{%
  \stepcounter{bigq}%
  \vspace{6pt}%
  \noindent{\large\bfseries 第\chinese{bigq}大题\quad #1}%
  \vspace{4pt}\par
}

\newcounter{subq}[bigq]
\newcommand{\SubQ}{%
  \stepcounter{subq}%
  \noindent\textbf{（\arabic{subq}）}\quad
}

% 分隔线
\newcommand{\examrule}{\noindent\rule{\linewidth}{0.6pt}}

% 单选题选项格式（缩进+紧凑排列）
\newcommand{\opts}[4]{%
  \par\noindent\vspace{2pt}%
  \hspace*{2em}\begin{minipage}[t]{\dimexpr\linewidth-2em\relax}
    \setlength{\parskip}{2pt}
    A.\ #1\par
    B.\ #2\par
    C.\ #3\par
    D.\ #4
  \end{minipage}\par\vspace{4pt}
}

% 多选题选项格式
\newcommand{\mopts}[4]{%
  \par\noindent\vspace{2pt}%
  \hspace*{2em}\begin{minipage}[t]{\dimexpr\linewidth-2em\relax}
    \setlength{\parskip}{2pt}
    A.\ #1\par
    B.\ #2\par
    C.\ #3\par
    D.\ #4
  \end{minipage}\par\vspace{4pt}
}

% ---- 文档开始 ----
\begin{document}

% ============================================================
%  封面
% ============================================================
\begin{center}
  {\LARGE\bfseries 习近平新时代中国特色社会主义思想概论}\\[8pt]
  {\large\bfseries 全真模拟仿真试卷（第二套·含参考答案）}\\[16pt]
  \begin{tabular}{lll}
    考试时间：120 分钟 & \quad & 满分：100 分 \\[4pt]
    \multicolumn{3}{l}{注意事项：请将答案写在答题纸指定区域，写在本卷上无效。}
  \end{tabular}
  \vspace{4pt}

  \examrule

  \vspace{6pt}
  \begin{tabular}{|c|c|c|c|c|c|}
    \hline
    \textbf{题号} & 一 & 二 & 三 & \textbf{总分} & \textbf{复核人} \\
    \hline
    \textbf{满分} & 30 & 20 & 50 & 100 & \\
    \hline
    \textbf{得分} & & & & & \\
    \hline
  \end{tabular}
\end{center}

\vspace{10pt}
\examrule
\vspace{6pt}

% ============================================================
%  第一大题：单项选择题（30 题，每题 1 分，共 30 分）
% ============================================================
\BigQ{单项选择题（每题 1 分，共 30 分。每题只有一个正确选项）}

\SubQ 党的十八大以来，党和国家事业取得历史性成就、发生历史性变革，最根本的原因在于有习近平总书记作为党中央的核心、全党的核心掌舵领航，在于有（\quad）科学指引。
\opts{邓小平理论}{"三个代表"重要思想}{科学发展观}{习近平新时代中国特色社会主义思想}

\SubQ "两个结合"重大论断中，"根脉"指的是（\quad）。
\opts{马克思主义基本原理}{中华优秀传统文化}{中国革命文化}{社会主义先进文化}

\SubQ "两个结合"重大论断中，"魂脉"指的是（\quad）。
\opts{马克思主义基本原理}{中华优秀传统文化}{中国革命文化}{世界优秀文明成果}

\SubQ 中国特色社会主义是党和人民历经千辛万苦、付出巨大代价取得的（\quad）。
\opts{重要成就}{根本成就}{基础成就}{关键成就}

\SubQ 中国特色社会主义是科学社会主义理论逻辑和（\quad）的辩证统一。
\opts{中国社会发展历史逻辑}{世界社会主义发展逻辑}{国际共产主义运动逻辑}{马克思主义发展逻辑}

\SubQ 一个国家实行什么主义，关键要看这个主义能否（\quad）。
\opts{符合本国国情}{解决这个国家面临的历史性课题}{顺应世界潮流}{得到人民拥护}

\SubQ 我国发展仍处于（\quad）的基本国情没有变。
\opts{社会主义初级阶段}{社会主义中级阶段}{社会主义高级阶段}{社会主义过渡阶段}

\SubQ 我国是世界最大发展中国家的（\quad）没有变。
\opts{国际地位}{国际影响}{综合国力}{发展阶段}

\SubQ "两个奇迹"指的是（\quad）。
\opts{经济快速发展奇迹和社会长期稳定奇迹}{经济快速发展奇迹和科技跨越发展奇迹}{社会长期稳定奇迹和疫情防控奇迹}{脱贫攻坚奇迹和抗击疫情奇迹}

\SubQ 党的基本路线是国家的（\quad）。
\opts{生命线、人民的幸福线}{政治线、人民的利益线}{发展线、人民的富裕线}{根本线、人民的保障线}

\SubQ 党的基本路线可以概括为（\quad）。
\opts{一个中心、两个基本点}{两个中心、一个基本点}{一个中心、三个基本点}{三个中心、两个基本点}

\SubQ 中国梦的本质是国家富强、民族振兴、（\quad）。
\opts{人民幸福}{人民富裕}{人民安康}{人民富足}

\SubQ 全面建成社会主义现代化强国的战略安排中，从2035年到本世纪中叶，要把我国建成（\quad）。
\opts{社会主义现代化强国}{基本实现现代化的国家}{中等发达国家}{全面建成小康社会的国家}

\SubQ "四个全面"战略布局中，居于引领地位的是（\quad）。
\opts{全面建设社会主义现代化国家}{全面深化改革}{全面依法治国}{全面从严治党}

\SubQ 中国式现代化的本质要求中，根本保证是（\quad）。
\opts{坚持中国共产党领导}{坚持中国特色社会主义}{实现高质量发展}{发展全过程人民民主}

\SubQ 中国式现代化是（\quad）的现代化。
\opts{人口规模巨大}{人口数量稳定}{人口结构优化}{人口质量提升}

\SubQ 团结奋斗是中国共产党、中国人民最显著的（\quad）。
\opts{精神标识}{政治标识}{文化标识}{历史标识}

\SubQ 中国共产党是（\quad）。
\opts{最高政治领导力量}{最高行政领导力量}{最高立法领导力量}{最高监察领导力量}

\SubQ 维护党中央权威和集中统一领导，是（\quad）的重大建党原则。
\opts{一个成熟的马克思主义执政党}{一个先进的马克思主义革命党}{一个奋斗的马克思主义执政党}{一个开放的马克思主义执政党}

\SubQ 党的领导是（\quad）。
\opts{全面的、系统的、整体的}{局部的、系统的、整体的}{全面的、分散的、整体的}{全面的、系统的、部分的}

\SubQ 共同富裕是中国式现代化的（\quad）。
\opts{重要特征}{基本特征}{一般特征}{次要特征}

\SubQ 共同富裕要坚持（\quad）推进。
\opts{有底线、渐进式}{齐头并进式}{同步达到式}{分批集中式}

\SubQ 新时代全面深化改革开放，就其艰巨性、复杂性和系统性来说，是一场（\quad）。
\opts{深刻革命}{深刻变革}{深刻调整}{深刻转型}

\SubQ 全面深化改革必须以（\quad）为主线。
\opts{制度建设}{经济建设}{政治建设}{文化建设}

\SubQ 全面深化改革要以（\quad）为重点。
\opts{经济体制改革}{政治体制改革}{文化体制改革}{社会体制改革}

\SubQ 改革和法治如（\quad）。
\opts{鸟之两翼、车之两轮}{舟之两桨、马之四蹄}{山之两峰、河之两岸}{树之两枝、楼之两柱}

\SubQ 在新发展理念中，注重解决发展不平衡问题的是（\quad）。
\opts{协调}{创新}{绿色}{开放}

\SubQ 在新发展理念中，注重解决社会公平正义问题的是（\quad）。
\opts{共享}{创新}{协调}{绿色}

\SubQ 高质量发展，就是从（\quad）转向"好不好"。
\opts{"有没有"}{"多不多"}{"快不快"}{"大不大"}

\SubQ 构建新发展格局最本质的特征是（\quad）。
\opts{实现高水平的自立自强}{扩大内需}{深化供给侧结构性改革}{推动产业链供应链优化升级}

\vspace{10pt}
\examrule

% ============================================================
%  第二大题：多项选择题（10 题，每题 2 分，共 20 分）
% ============================================================
\BigQ{多项选择题（每题 2 分，共 20 分。每题有两个或两个以上正确选项，少选、多选、错选均不得分）}

\SubQ 习近平新时代中国特色社会主义思想的历史地位体现在（\quad）。
\mopts{当代中国马克思主义、二十一世纪马克思主义}{中华文化和中国精神的时代精华}{实现了马克思主义中国化时代化新的飞跃}{是党和国家必须长期坚持的指导思想}

\SubQ 习近平新时代中国特色社会主义思想是完整的科学体系，其构成包括（\quad）。
\mopts{十个明确}{十四个坚持}{十三个方面成就}{六个必须坚持}

\SubQ "六个必须坚持"包括（\quad）。
\mopts{必须坚持人民至上}{必须坚持自信自立}{必须坚持问题导向}{必须坚持系统观念}

\SubQ 中国式现代化的本质要求包括（\quad）。
\mopts{坚持中国共产党领导、坚持中国特色社会主义}{实现高质量发展、发展全过程人民民主}{丰富人民精神世界、实现全体人民共同富裕}{促进人与自然和谐共生、推动构建人类命运共同体、创造人类文明新形态}

\SubQ 推进中国式现代化需要正确处理的重大关系包括（\quad）。
\mopts{顶层设计与实践探索的关系}{战略与策略的关系}{守正与创新的关系}{效率与公平的关系}

\SubQ 中国共产党领导是中国特色社会主义制度的最大优势，主要体现在（\quad）。
\mopts{党的领导制度是我国的根本领导制度}{党的领导是做好党和国家各项工作的根本保证}{党的领导是战胜一切困难和风险的"定海神针"}{党的自身优势是中国特色社会主义制度优势的主要来源}

\SubQ 推动全体人民共同富裕取得更为明显的实质性进展，要处理好（\quad）的关系。
\mopts{先富和共富}{效率和公平}{发展和共享}{局部和全局}

\SubQ 全面深化改革开放要坚持的正确方向包括（\quad）。
\mopts{必须坚持和改善党的全面领导、坚持和完善中国特色社会主义制度}{必须坚持以人民为中心，促进社会公平正义、增进人民福祉}{必须有利于进一步解放思想}{必须有利于进一步解放和发展社会生产力、解放和增强社会活力}

\SubQ 全过程人民民主是（\quad）的民主。
\mopts{全链条}{全方位}{全覆盖}{最广泛、最真实、最管用}

\SubQ 中国特色社会主义政治制度的组成包括（\quad）。
\mopts{根本政治制度}{基本政治制度}{重要政治制度}{辅助政治制度}

\vspace{10pt}
\examrule

% ============================================================
%  第三大题：简答题（5 题，每题 10 分，共 50 分）
% ============================================================
\BigQ{简答题（每题 10 分，共 50 分。要求：在给出正确答案的基础上进行简要论述，否则不能得满分）}

\SubQ（10 分）试述中国创造人类文明新形态的三个方面。

\vspace{8pt}
\SubQ（10 分）试述人民立场是中国共产党的根本政治立场。

\vspace{8pt}
\SubQ（10 分）试述新时代全面深化改革开放是一场深刻革命及其原因。

\vspace{8pt}
\SubQ（10 分）试述全面深化改革的总目标，并说明如何统筹推进各领域各方面改革开放。

\vspace{8pt}
\SubQ（10 分）试述人民民主是社会主义的生命，并说明坚定不移走中国特色社会主义政治发展道路的根本要求。

\vspace{16pt}
\examrule
\vspace{4pt}
\begin{center}
  {\large\bfseries ——试题到此结束，请认真检查——}
\end{center}
\vspace{4pt}
\examrule

% ============================================================
%  参考答案
% ============================================================
\newpage
\setcounter{bigq}{0}

\begin{center}
  {\LARGE\bfseries 参考答案与评分标准}
\end{center}
\vspace{8pt}
\examrule

% -------- 第一大题参考答案 --------
\BigQ{单项选择题参考答案}

\begin{answerbox}
\renewcommand{\arraystretch}{1.4}
\begin{tabularx}{\linewidth}{@{}lX@{}}
1.\ D & 2.\ B \quad 3.\ A \quad 4.\ B \quad 5.\ A \quad 6.\ B \\
7.\ A & 8.\ A \quad 9.\ A \quad 10.\ A \quad 11.\ A \quad 12.\ A \\
13.\ A & 14.\ A \quad 15.\ A \quad 16.\ A \quad 17.\ A \quad 18.\ A \\
19.\ A & 20.\ A \quad 21.\ A \quad 22.\ A \quad 23.\ A \quad 24.\ A \\
25.\ A & 26.\ A \quad 27.\ A \quad 28.\ A \quad 29.\ A \quad 30.\ A \\
\end{tabularx}

\medskip
\noindent\textbf{评分说明}：每题 1 分，共 30 分。所选答案与参考答案完全一致方可得分，错选、多选、不选均不得分。
\end{answerbox}

% -------- 第二大题参考答案 --------
\BigQ{多项选择题参考答案}

\begin{answerbox}
\renewcommand{\arraystretch}{1.4}
\begin{tabularx}{\linewidth}{@{}lX@{}}
1.\ ABCD & 2.\ ABCD \quad 3.\ ABCD \quad 4.\ ABCD \quad 5.\ ABCD \\
6.\ ABCD & 7.\ ABC \quad 8.\ ABCD \quad 9.\ ABC \quad 10.\ ABC \\
\end{tabularx}

\medskip
\noindent\textbf{评分说明}：每题 2 分，共 20 分。每题有两个或两个以上正确选项，少选、多选、错选均不得分。
\end{answerbox}

% -------- 第三大题参考答案 --------
\BigQ{简答题参考答案}

\begin{answerbox}
\textbf{（1）试述中国创造人类文明新形态的三个方面。（10 分）}

\textbf{【参考答案】}

中国式现代化创造了人类文明新形态，这一新形态的创造主要体现在以下三个方面：

\textbf{第一，提供了全新的现代化路径，超越了西方现代化模式。} 中国式现代化深深植根于中华优秀传统文化，体现科学社会主义的先进本质，借鉴吸收一切人类优秀文明成果，代表人类文明进步的发展方向。它打破了"现代化=西方化"的迷思，展现了不同于西方现代化模式的新图景。西方现代化以资本为中心、遵循两极分化、物质主义膨胀、对外扩张掠夺的逻辑；而中国式现代化则坚持人民至上、追求共同富裕、促进物质精神共同进步、走和平发展道路，为人类文明发展提供了崭新范式。

\textbf{第二，为发展中国家提供了新的选择。} 长期以来，发展中国家在追求现代化的进程中面临着"依附"还是"脱钩"的两难困境。中国式现代化的成功实践证明，每个国家完全可以基于自身历史文化传统和现实国情，走出一条符合自身特点的现代化道路。广大发展中国家可以借鉴中国式现代化的经验，独立自主地探索适合本国国情的现代化路径，不必照搬西方模式，这为人类文明多样性发展开辟了广阔空间。

\textbf{第三，需牢牢把握推进中国式现代化的重大原则。} 创造人类文明新形态，必须牢牢把握坚持和加强党的全面领导、坚持中国特色社会主义道路、坚持以人民为中心的发展思想、坚持深化改革开放、坚持发扬斗争精神五大重大原则。这五大原则与"四项基本原则"一脉相承又与时俱进，共同构成推进中国式现代化、创造人类文明新形态的根本政治保证、根本方向指引、根本价值取向、根本动力源泉和根本精神支撑。

\textbf{【评分标准】} 三个方面每个方面 3 分，共 9 分；综合论述加 1 分，共 10 分。要点准确、论述充分方可得满分。
\end{answerbox}

\begin{answerbox}
\textbf{（2）试述人民立场是中国共产党的根本政治立场。（10 分）}

\textbf{【参考答案】}

人民立场是中国共产党的根本政治立场，这一论断具有深刻的内涵和现实意义：

\textbf{第一，人民是历史的创造者，是真正的英雄。} 马克思主义唯物史观认为，人民群众是历史的创造者，是推动社会发展的决定性力量。中国共产党始终把人民作为创造历史的真正动力，作为历史发展和社会进步的主体力量。党的一切奋斗都是为了人民，党的一切成就都依靠人民创造。这一立场是党区别于其他政党的根本标志，是党之所以能够赢得人民衷心拥护的根本原因。

\textbf{第二，民心是最大的政治，决定事业兴衰成败。} "江山就是人民，人民就是江山"，打江山、守江山，守的是人民的心。民心向背关系党的生死存亡，关系事业兴衰成败。只有紧紧依靠最广大人民群众，为人民执好政、掌好权，才能始终赢得人民群众衷心拥护，实现长期执政，带领人民创造更加美好的生活。这一论断深刻揭示了人心向背的政治分量，是党治国理政的根本遵循。

\textbf{第三，人民立场是根本政治立场的具体体现。} 坚持人民立场，就要始终牢记党的初心和使命，即为中国人民谋幸福、为中华民族谋复兴；就要始终保持党同人民群众的血肉联系，倾听群众呼声、回应群众期待；就要热爱人民、尊重人民、敬畏人民，把人民放在心中最高位置。人民性是马克思主义的本质属性和鲜明品格，坚持以人民为中心，是新时代坚持和发展中国特色社会主义的根本立场。

\textbf{第四，坚持人民立场的实践要求。} 要坚持人民至上，把人民对美好生活的向往作为奋斗目标；要依靠人民创造历史伟业，尊重人民主体地位和首创精神；要把人民作为党的工作的最高裁决者和最终评判者，让群众满意成为党做好一切工作的价值取向和根本标准；要全面落实以人民为中心的发展思想，推动全体人民共同富裕取得更为明显的实质性进展。

\textbf{【评分标准】} 四个方面每个方面 2.5 分，共 10 分。要点准确、论述充分方可得满分；如能联系"江山就是人民"等论断加以深化，可酌情加分（但总分不超过 10 分）。
\end{answerbox}

\begin{answerbox}
\textbf{（3）试述新时代全面深化改革开放是一场深刻革命及其原因。（10 分）}

\textbf{【参考答案】}

\textbf{一、基本判断（2 分）}

新时代全面深化改革开放，就其艰巨性、复杂性和系统性来说，是一场深刻的革命。这一判断深刻揭示了新时代改革开放的本质特征和重大意义，是党中央对全面深化改革开放作出的重大政治判断和理论概括。

\textbf{二、深刻革命的原因（8 分）}

\textbf{第一，全面深化改革开放是涉险滩、闯难关的变革。} 经过四十多年的改革开放，容易改的、好吃的小肉都吃掉了，剩下的都是难啃的硬骨头。新时代改革开放进入了攻坚期和深水区，涉及深层次利益格局调整，触及深层次思想观念束缚，破解发展面临的新的难题、化解来自各方面的新的风险挑战，必须进行具有许多新的历史特点的伟大斗争。

\textbf{第二，全面深化改革开放是一场思想理论的深刻变革。} 新时代改革开放冲破了因循守旧、固步自封的思想束缚，破除了体制机制弊端和利益固化藩篱，实现了思想理论的重大飞跃。它以全新的视野深化了对共产党执政规律、社会主义建设规律、人类社会发展规律的认识，是新时代思想解放的集中体现。

\textbf{第三，全面深化改革开放是一场改革组织方式的深刻变革。} 党中央加强对全面深化改革的顶层设计、总体规划、统筹协调，建立了更加科学高效的改革领导机制、决策机制、协调机制、推进机制、督察落实机制。改革的系统性、整体性、协同性显著增强，整体推进和重点突破相结合成为基本方法。

\textbf{第四，全面深化改革开放是一场国家制度和治理体系的深刻变革。} 党的十八届三中全会首次提出全面深化改革总目标，即完善和发展中国特色社会主义制度、推进国家治理体系和治理能力现代化。各领域基础性制度框架基本确立，许多领域实现历史性变革、系统性重塑、整体性重构，国家治理体系和治理能力现代化水平明显提高。

\textbf{第五，全面深化改革开放是一场人民广泛参与的深刻变革。} 改革为了人民、改革依靠人民、改革成果由人民共享。新时代改革开放充分调动了人民群众的积极性、主动性、创造性，形成了改革开放的强大合力，是党领导人民进行的伟大社会革命的继续。

\textbf{【评分标准】} 基本判断 2 分；五个原因每个方面 1.5 分，共 7.5 分；综合提升 0.5 分，共 10 分。要点准确、论述充分方可得满分。
\end{answerbox}

\begin{answerbox}
\textbf{（4）试述全面深化改革的总目标，并说明如何统筹推进各领域各方面改革开放。（10 分）}

\textbf{【参考答案】}

\textbf{一、全面深化改革的总目标（4 分）}

全面深化改革的总目标是：完善和发展中国特色社会主义制度、推进国家治理体系和治理能力现代化。

\textbf{第一，"完善和发展中国特色社会主义制度"规定了改革的根本方向。} 无论改什么、怎么改，都要坚持中国共产党领导、坚持中国特色社会主义，就是要通过改革推动中国特色社会主义制度更加成熟更加定型、更好发挥中国特色社会主义制度的优越性。这一句话强调的是"方向"，确保改革始终沿着正确方向前进，既不走封闭僵化的老路，也不走改旗易帜的邪路。

\textbf{第二，"推进国家治理体系和治理能力现代化"明确了改革的鲜明指向和时代要求。} 就是要通过改革进一步增强我国制度活力，把制度优势转化为国家治理效能。国家治理体系和治理能力是一个有机整体，相辅相成，有了好的国家治理体系才能提高国家治理能力，提高国家治理能力才能充分发挥国家治理体系的效能。

\textbf{第三，两句话构成一个有机整体。} 这一总目标是一个内涵丰富的有机整体，两句话都讲，才是完整的、全面的，不能割裂开来理解。

\textbf{二、如何统筹推进各领域各方面改革开放（6 分）}

统筹推进各领域各方面改革开放，必须以全面深化改革总目标为引领，重点把握以下四个方面：

\textbf{第一，以制度建设为主线。} 制度建设贯穿改革全过程，要抓住制度建设这个主线，把解决实际问题与制度建设结合起来，把改革成果用制度形式固定下来。要通过改革，构建系统完备、科学规范、运行有效的制度体系，使各方面制度更加成熟更加定型。

\textbf{第二，以经济体制改革为重点。} 经济基础决定上层建筑，经济体制改革对其他方面改革具有重要影响和传导作用。要紧紧围绕使市场在资源配置中起决定性作用和更好发挥政府作用，深化经济体制改革，牵引和带动其他领域改革协同推进。

\textbf{第三，以重要领域和关键环节改革为突破口。} 改革不能平均用力、齐头并进，要注重抓主要矛盾和矛盾的主要方面，注重抓重要领域和关键环节。通过重点突破带动全局，以重点领域和关键环节的突破带动改革全局推进。

\textbf{第四，实施好"六个紧紧围绕"的路线图。} 紧紧围绕使市场在资源配置中起决定性作用深化经济体制改革；紧紧围绕坚持党的领导、人民当家作主、依法治国有机统一深化政治体制改革；紧紧围绕建设社会主义核心价值体系深化文化体制改革；紧紧围绕更好保障和改善民生深化社会体制改革；紧紧围绕建设美丽中国深化生态文明体制改革；紧紧围绕提高科学执政、民主执政、依法执政水平深化党的建设制度改革。使各方面改革协同推进、环环相扣、形成合力。

\textbf{【评分标准】} 总目标 4 分（两句话各 2 分）；统筹推进四个方面每个方面 1.5 分，共 6 分。要点准确、论述充分方可得满分。
\end{answerbox}

\begin{answerbox}
\textbf{（5）试述人民民主是社会主义的生命，并说明坚定不移走中国特色社会主义政治发展道路的根本要求。（10 分）}

\textbf{【参考答案】}

\textbf{一、人民民主是社会主义的生命（5 分）}

\textbf{第一，民主是全人类共同价值，是中国共产党和中国人民始终不渝坚持的重要理念。} 民主是历史的、具体的，一个国家实行什么样的民主制度，走什么样的民主发展道路，必须与这个国家的国情相适应。实现民主的道路并非只有一条，社会性质不同、现实国情不同，则民主道路不同、民主形态各异。

\textbf{第二，人民民主建立在社会主义经济基础之上，体现了社会主义国家的性质。} 人民民主反映了社会主义制度的本质要求，是一种新型的社会主义民主。它保证了人民当家作主，是社会主义优越性的重要体现。没有民主就没有社会主义，就没有社会主义现代化，人民民主是社会主义的生命这一论断，深刻揭示了民主与社会主义的内在统一关系。

\textbf{第三，人民民主是全面建设社会主义现代化国家的应有之义。} 发展社会主义民主政治，是全面建设社会主义现代化国家的内在要求和重要目标。社会主义现代化是全体人民的现代化，必然要求具有高度发达的民主政治，以实现好、维护好、发展好最广大人民根本利益为出发点和落脚点，保障人民当家作主。全面建成社会主义现代化强国，人民是决定性力量。

\textbf{第四，全过程人民民主是社会主义民主政治的本质属性。} 全过程人民民主是过程民主和成果民主的统一、程序民主和实质民主的统一、直接民主和间接民主的统一、人民民主和国家意志的统一，是最广泛、最真实、最管用的民主。注意区分"成果民主"与"结果民主"，我们强调的是"成果民主"，不是"结果民主"。

\textbf{二、坚定不移走中国特色社会主义政治发展道路的根本要求（5 分）}

坚定不移走中国特色社会主义政治发展道路，必须做到以下三点：

\textbf{第一，必须坚持党的领导、人民当家作主、依法治国有机统一。} 党的领导是人民当家作主和依法治国的根本保证，人民当家作主是社会主义民主政治的本质特征，依法治国是党领导人民治理国家的基本方式。三者统一于我国社会主义民主政治伟大实践。走中国特色社会主义政治发展道路，最根本的是坚持党的领导，这是政治发展道路的核心要义。

\textbf{第二，必须积极稳妥推进政治体制改革。} 政治体制改革是我国全面深化改革的重要组成部分，必须坚持正确方向，坚持中国特色社会主义政治制度自信，以是否有利于坚持和改善党的全面领导、坚持和完善中国特色社会主义制度为根本标准，积极稳妥加以推进。要不断健全人民当家作主制度体系，发展全过程人民民主，丰富民主形式，拓宽民主渠道。

\textbf{第三，必须始终保持政治定力。} 要坚定走中国特色社会主义政治发展道路的信心和决心，不能照搬他国政治制度模式，不能放弃我国社会主义政治制度的根本。要保持清醒头脑，增强政治敏锐性和政治鉴别力，坚决抵制西方"宪政民主"、多党轮流执政、"三权鼎立"等政治模式的影响，确保我国政治发展道路始终沿着正确方向前进。

\textbf{【评分标准】} "人民民主是社会主义的生命"四个方面共 5 分；"走中国特色社会主义政治发展道路"三个根本要求每个方面 1.5 分，共 4.5 分；综合提升 0.5 分，共 10 分。要点准确、论述充分方可得满分。
\end{answerbox}

\vspace{12pt}
\examrule
\begin{center}
  \textit{参考答案仅供参考，评分时应酌情给分，步骤正确可得过程分。}
\end{center}

\end{document}
